\documentclass[14pt]{extarticle}
\usepackage{preamble}
\geometry{letterpaper, landscape, margin=0.5in}
\renewcommand{\answer}[1]{}
\setlength{\parskip}{4pt}

\begin{document}
\pagestyle{empty}

\daybanner{Unit 5 \textperiodcentered\ Topic 5.1 \textperiodcentered\ THE PROBLEM}{One Plot, Two Reporting Strategies}

\noindent\begin{minipage}[t]{0.57\linewidth}
\vspace{0pt}
Suppose a test kitchen measures dough volume for ten batches at assigned times after mixing. The scatterplot shows the paired measurements.\par\smallskip \textbf{Jules:} Use a straight-line summary because volume generally increases with time. Treat the 50-minute batch as ordinary variation.\par \textbf{Ravi:} Decide the form from how the rate changes, and report the 50-minute batch separately if it departs from nearby values.\par Selected measurements (minutes, mL): $(0,420)$, $(20,590)$, $(70,770)$, $(90,785)$.\par
\end{minipage}\hfill
\begin{minipage}[t]{0.40\linewidth}
\vspace{0pt}
\begin{center}
\begin{tikzpicture}
\begin{axis}[
  width=0.70\linewidth,
  height=1.89in,
  xmin=0, xmax=90, ymin=380, ymax=820,
  xlabel={Minutes after mixing}, ylabel={Dough volume (mL)}, grid=major,
  tick label style={font=\small}, label style={font=\small}
]
\addplot[only marks, mark=*, mark size=1.8pt, color=dayblue] coordinates {
  (0,420)
  (10,510)
  (20,590)
  (30,650)
  (40,695)
  (50,600)
  (60,750)
  (70,770)
  (80,780)
  (90,785)
};
\end{axis}
\end{tikzpicture}
\end{center}
\end{minipage}

\begin{firsttakebox}{FIRST TAKE --- one sentence before you work the parts}
Which reporting strategy seems more defensible, Jules's or Ravi's? Cite one graph feature.
\end{firsttakebox}

\begin{description}[leftmargin=1.15in, labelwidth=1in, itemsep=2pt, topsep=2pt]
  \item[\dokbadge{1}~(a)] Name the explanatory and response variables, with units.
  \item[\dokbadge{2}~(b)] Find the volume increase from 0 to 20 minutes and from 70 to 90 minutes. Is the 50-minute batch above or below the nearby batches?
  \item[\dokbadge{3}~(c)] Whose reporting strategy is supported? Describe the direction, form, strength, and unusual point in context, using your two increases. Explain how a straight-line summary could mislead someone predicting later growth. Write 3--4 sentences.
\end{description}

\vfill
\textbf{Self-paced bonus problem. Work (a)--(c) from this sheet; turn it in whenever you finish.}

\end{document}
